\documentclass[11pt]{article}
\usepackage{mathblatt}
\def\mitzone{1}
\begin{document}
\blattkopf{Winkel und Dreiecke}{Gesamt}
\weit
\verzeichniszeile{\verz{zone}{Kennst du schon (Nr.~1–7)} \verztrenn \verz{e1}{1 Winkel messen und zeichnen (Nr.~8–11)} \verztrenn \verz{e2}{2 Winkel an Geradenkreuzungen und Parallelen (Nr.~12–15)} \verztrenn \verz{e3}{3 Winkelsummen, Dreiecke und Vierecke (Nr.~16–19)} \verztrenn \verz{e4}{4 Dreiecke konstruieren (Nr.~20–22)} \verztrenn \verz{e5}{5 Besondere Linien im Dreieck und Satz des Thales (Nr.~23–24)} \verztrenn \verz{abhaken}{Das kann ich}}
% Zone „Das kennst du schon“ – aus bank/winkel-dreiecke/zone.jsonl (zusammenbau v0.1)
\einheitenkopf[zone][Kennst du schon]{Das kennst du schon}
\setcounter{aufgabe}{0}

%% TODO Titel: Ich-kann-Satz fehlt in der Bank (Platzhalter: Kettenname) – Nr. 1
%% TODO Anweisung: ein Satz über der ersten Teilaufgabe fehlt in der Bank – Nr. 1
\begin{aufgabe}{Mit Lineal und Geodreieck zeichnen}
\begin{teile}
\teil Miss mit dem Lineal die Länge und die Breite des Rechtecks.

\rechteck{5}{3}
Länge: \leerfeld[cm] , Breite: \leerfeld[cm]
\teil Zeichne eine Gerade g und einen Punkt P, der 3 cm von g entfernt liegt. Zeichne mit dem Geodreieck die Senkrechte zu g durch P und die Parallele zu g durch P.

\begin{ksys}[xmin=0,xmax=10,ymin=0,ymax=6] \end{ksys}
\end{teile}
\end{aufgabe}

%% TODO Titel: Ich-kann-Satz fehlt in der Bank (Platzhalter: Kettenname) – Nr. 2
%% TODO Anweisung: ein Satz über der ersten Teilaufgabe fehlt in der Bank – Nr. 2
\begin{aufgabe}{Addieren und Subtrahieren im Bereich bis zum Vollwinkel, auch mit Dezimalzahlen (Vielfaches minus zwei Dezimalgrade)}
\begin{teile}
\teil $38^\circ + 29^\circ$ \leerfeld °
\teil $96{,}4^\circ + 38{,}7^\circ$ \leerfeld °
\end{teile}
\end{aufgabe}

%% TODO Titel: Ich-kann-Satz fehlt in der Bank (Platzhalter: Kettenname) – Nr. 3
%% TODO Anweisung: ein Satz über der ersten Teilaufgabe fehlt in der Bank – Nr. 3
\begin{aufgabe}{Längeneinheiten m und km umrechnen (1,5 km = 1500 m)}
\begin{teile}
\teil 3 km in Meter \leerfeld[m]
\teil 2,4 km in Meter \leerfeld[m]
\end{teile}
\end{aufgabe}

%% TODO Titel: Ich-kann-Satz fehlt in der Bank (Platzhalter: Kettenname) – Nr. 4
%% TODO Anweisung: ein Satz über der ersten Teilaufgabe fehlt in der Bank – Nr. 4
\begin{aufgabe}{Vierecksarten kennen}
\begin{teile}
\teil Wie heißt ein Viereck mit vier rechten Winkeln und vier gleich langen Seiten? \leerfeld
\teil Welche Seiten des Vierecks sind parallel? Wie heißt das Viereck?

\trapez{5}{3}{2.5}
\end{teile}
\end{aufgabe}

%% TODO Titel: Ich-kann-Satz fehlt in der Bank (Platzhalter: Kettenname) – Nr. 5
%% TODO Anweisung: ein Satz über der ersten Teilaufgabe fehlt in der Bank – Nr. 5
\begin{aufgabe}{Kreis mit dem Zirkel zeichnen, Radius und Durchmesser unterscheiden}
\begin{teile}
\teil Zeichne mit dem Zirkel einen Kreis mit dem Radius 3 cm.

\begin{ksys}[xmin=0,xmax=10,ymin=0,ymax=6] \end{ksys}
\teil Welche Strecke ist der Radius r, welche der Durchmesser d? Miss beide.

\begin{kreis}[2] \mittelpunkt{M} \radius{40}{r} \durchmesser{0}{d} \end{kreis}
r = \leerfeld[cm] , d = \leerfeld[cm]
\end{teile}
\end{aufgabe}

%% TODO Titel: Ich-kann-Satz fehlt in der Bank (Platzhalter: Kettenname) – Nr. 6
%% TODO Anweisung: ein Satz über der ersten Teilaufgabe fehlt in der Bank – Nr. 6
\begin{aufgabe}{Addieren und Subtrahieren im Bereich bis zum Vollwinkel, auch mit Dezimalzahlen (Vielfaches minus zwei Dezimalgrade) – Fehler finden}
\begin{teile}
\teil Mia rechnet: \rechnung{180^\circ - 46{,}8^\circ - 37{,}5^\circ &= 96{,}3^\circ} Finde den Fehler und rechne richtig.
\end{teile}
\end{aufgabe}

%% TODO Titel: Ich-kann-Satz fehlt in der Bank (Platzhalter: Kettenname) – Nr. 7
%% TODO Anweisung: ein Satz über der ersten Teilaufgabe fehlt in der Bank – Nr. 7
\begin{aufgabe}{Addieren und Subtrahieren im Bereich bis zum Vollwinkel, auch mit Dezimalzahlen (Vielfaches minus zwei Dezimalgrade) – selbst rechnen}
\begin{teile}
\teil $180^\circ - 38{,}6^\circ - 55{,}7^\circ$ \leerfeld °
\end{teile}
\end{aufgabe}

\clearpage
% Einheit 1 von 5 – Winkel messen und zeichnen (bank/winkel-dreiecke/e1.jsonl, zusammenbau v0.1)
%% TODO Zweigzeile Teil 1: was hier gelernt wird (Satz in Schülersprache aus dem Einheitstitel) – Einheit 1
\einheitenkopf[e1]{Einheit 1 von 5 · Winkel messen und zeichnen}
\zweigzeile{ab Kl. 5 · P10}
\setcounter{aufgabe}{7}

%% TODO Titel: Ich-kann-Satz fehlt in der Bank (Platzhalter: Kettenname) – Nr. 8
%% TODO Anweisung: ein Satz über der ersten Teilaufgabe fehlt in der Bank – Nr. 8
\begin{aufgabe}{Welche Skala?}
\begin{teile}
\teil Nicht ablesen: Gilt beim Messen von α am Geodreieck die kleinere oder die größere der beiden Zahlen? \\ \kreuz{die kleinere Zahl} \\ \kreuz{die größere Zahl}

\winkel[5]{38}{\alpha}
\end{teile}
\end{aufgabe}

%% TODO Titel: Ich-kann-Satz fehlt in der Bank (Platzhalter: Kettenname) – Nr. 9
%% TODO Anweisung: ein Satz über der ersten Teilaufgabe fehlt in der Bank – Nr. 9
\begin{aufgabe}{Winkel messen}
%% TODO \rechenplatz{n}: Schrittzahl des Grundfalls fehlt in der Bank (nur bei fester Schreibform, 2.3 b)
\begin{teile}
\teil Welche Winkelart hat α? Nicht messen. \\ \kreuz{spitz} \\ \kreuz{recht} \\ \kreuz{stumpf} \\ \kreuz{gestreckt} \\ \kreuz{überstumpf}

\winkel[4]{27}{\alpha}
\teil Miss den Winkel α mit dem Geodreieck.

\winkel[5]{37}{\alpha}
α = \leerfeld °
\teil Miss den stumpfen Winkel α mit dem Geodreieck.

\winkel[5]{118}{\alpha}
α = \leerfeld °
\teil Zeichne an den Strahl mit dem Scheitel S einen Winkel von $55^\circ$.

\winkelstrahl
\teil Bestimme den überstumpfen Winkel α. Miss dazu den Rest bis zum Vollwinkel.

\winkel[4]{215}{\alpha}
α = \leerfeld °
\teil Schreibe den Winkel α im Dreieck ABC mit drei Buchstaben ($\angle$…). Welcher Punkt ist der Scheitel?

\dreieck{(0,0)}{(5,0)}{(1.5,3)}{a}{b}{c}{\alpha}{\beta}{\gamma}
\teil Die Winkel $34^\circ$ und $47^\circ$ liegen am selben Scheitel nebeneinander. Wie groß ist der Winkel, den beide zusammen bilden? \leerfeld °
\teil Ein Winkel von $96^\circ$ ist durch einen Strahl in zwei Teilwinkel geteilt. Der eine misst $38^\circ$. Wie groß ist der andere? \leerfeld °
\teil Ein Wanderweg ist 1,2 km lang. Nach 450 m macht Ella eine Pause. Wie viele Meter hat sie noch vor sich? \leerfeld[m]
\teil Vom Punkt A aus sieht Jana die Spitze T eines Turms unter einem Winkel von $41^\circ$ zur waagerechten Grundlinie, die Spitze F der Fahne auf dem Turm unter $44^\circ$. Wie groß ist der Winkel γ zwischen den Sichtlinien AT und AF? \hfill (P10 2018 OS) \\ γ = \leerfeld °
\end{teile}
\end{aufgabe}

%% TODO Titel: Ich-kann-Satz fehlt in der Bank (Platzhalter: Kettenname) – Nr. 10
%% TODO Anweisung: ein Satz über der ersten Teilaufgabe fehlt in der Bank – Nr. 10
\begin{aufgabe}{Winkel schätzen}
\begin{teile}
\teil Schätze, ohne zu messen: Wie groß ist α im Vergleich zu einem rechten Winkel? \\ \kreuz{ein Viertel} \\ \kreuz{die Hälfte} \\ \kreuz{drei Viertel}

\winkel[5]{23}{\alpha}
\end{teile}
\end{aufgabe}

%% TODO Titel: Ich-kann-Satz fehlt in der Bank (Platzhalter: Kettenname) – Nr. 11
%% TODO Anweisung: ein Satz über der ersten Teilaufgabe fehlt in der Bank – Nr. 11
\begin{aufgabe}{Winkel messen – Fehler finden, Begründen, Anwendung}
\begin{teile}
\teil Tim hat den Winkel α gemessen und notiert: $\alpha = 144^\circ$. Finde den Fehler und gib den richtigen Wert an.

\winkel[5]{36}{\alpha}
\teil Warum heißt ein Winkel von $91^\circ$ stumpf, ein Winkel von $89^\circ$ aber spitz? Begründe.
\teil Eine Parkhausschranke dreht sich beim Öffnen aus der waagerechten Lage, bis sie senkrecht steht. Nach zwei Sekunden hat sie sich um $54^\circ$ gedreht. Um wie viel Grad muss sie sich noch drehen? \leerfeld °
\end{teile}
\end{aufgabe}

\clearpage
% Einheit 2 von 5 – Winkel an Geradenkreuzungen und Parallelen (bank/winkel-dreiecke/e2.jsonl, zusammenbau v0.1)
%% TODO Zweigzeile Teil 1: was hier gelernt wird (Satz in Schülersprache aus dem Einheitstitel) – Einheit 2
\einheitenkopf[e2]{Einheit 2 von 5 · Winkel an Geradenkreuzungen und Parallelen}
\zweigzeile{ab Kl. 6, je nach Buch bis Kl. 7 (am Gymnasium ab Kl. 5) · P10 oft}
\setcounter{aufgabe}{11}

%% TODO Titel: Ich-kann-Satz fehlt in der Bank (Platzhalter: Kettenname) – Nr. 12
%% TODO Anweisung: ein Satz über der ersten Teilaufgabe fehlt in der Bank – Nr. 12
\begin{aufgabe}{Sind die Geraden parallel?}
\begin{teile}
\teil Angabe: Die Geraden g und h sind parallel; die Gerade k schneidet beide. Darfst du g und h als parallel ansehen? Markiere Parallelen mit Pfeilen oder kreuze an: \\ \kreuz{als parallel angegeben} \\ \kreuz{nicht angegeben}

\parallelenpaar{58}{}{}{}{}
\end{teile}
\end{aufgabe}

%% TODO Titel: Ich-kann-Satz fehlt in der Bank (Platzhalter: Kettenname) – Nr. 13
%% TODO Anweisung: ein Satz über der ersten Teilaufgabe fehlt in der Bank – Nr. 13
\begin{aufgabe}{Winkel an Geraden}
%% TODO \rechenplatz{n}: Schrittzahl des Grundfalls fehlt in der Bank (nur bei fester Schreibform, 2.3 b)
\begin{teile}
\teil Wie liegen α und γ an der Kreuzung zueinander? \\ \kreuz{Scheitelwinkel} \\ \kreuz{Nebenwinkel} \\ \kreuz{Stufenwinkel} \\ \kreuz{Wechselwinkel}

\geradenkreuzung{55}{\alpha}{\beta}{\gamma}{\delta}
\teil An einer Geradenkreuzung ist $\alpha = 38^\circ$. Wie groß ist der Scheitelwinkel von α?

\geradenkreuzung{38}{\alpha}{}{}{}
\leerfeld °
\teil An einer Geradenkreuzung ist $\alpha = 44^\circ$. Wie groß ist ein Nebenwinkel von α?

\geradenkreuzung{44}{\alpha}{}{}{}
\leerfeld °
\teil An einer Geradenkreuzung ist $\alpha = 72^\circ$; β und δ liegen neben α, γ liegt α gegenüber. Gib β, γ und δ an.

\geradenkreuzung{72}{\alpha}{\beta}{\gamma}{\delta}
β = \leerfeld °, γ = \leerfeld °, δ = \leerfeld °
\teil Die Geraden g und h schneiden sich; zwischen ihnen liegt $\alpha = 64^\circ$. Eine dritte Gerade durch den Schnittpunkt teilt den Scheitelwinkel von α in zwei Teile, der untere misst $23^\circ$. Wie groß ist der obere Teil β? \hfill (P10 2014 OS) \\ β = \leerfeld °
\teil $g \parallel h$, g liegt oben, h unten, die Gerade k schneidet beide. An g liegt oberhalb von g rechts von k ein Winkel von $58^\circ$. Wie groß ist der Winkel an h oberhalb von h rechts von k?

\parallelenpaar{58}{}{}{}{}
\leerfeld °
\teil $g \parallel h$, g liegt oben, h unten, die Gerade k schneidet beide. An g liegt unterhalb von g links von k ein Winkel von $49^\circ$. Wie groß ist der Winkel an h oberhalb von h rechts von k? \hfill (P10 2020 OS) \\

\parallelenpaar{49}{}{}{}{}
\leerfeld °
\teil $g \parallel h$, g liegt oben, h unten, die Gerade k schneidet beide. An g liegt oberhalb von g rechts von k ein Winkel von $62^\circ$. Wie groß ist der Winkel an h oberhalb von h links von k?

\parallelenpaar{62}{}{}{}{}
\leerfeld °
\teil Im Dreieck ABC ist der Winkel bei C $58^\circ$ groß. Die Seiten AB und CB werden über B hinaus verlängert; zwischen den beiden Verlängerungen liegt ein Winkel von $43^\circ$. Wie groß ist der Innenwinkel β bei B? \hfill (P10 2019 OS) \\ β = \leerfeld °
\teil Im Parallelogramm ABCD ist $\alpha = 64^\circ$. Gib β, γ und δ an.

\parallelogramm[winkel={\alpha,\beta,\gamma,\delta}]{4}{2.5}{64}
β = \leerfeld °, γ = \leerfeld °, δ = \leerfeld °
\teil Die parallelen Geraden h (oben) und g (unten) werden von der Geraden k geschnitten. An h liegt unterhalb von h links von k ein Winkel von $44^\circ$; α liegt an g oberhalb von g links von k. Wie groß ist α? \hfill (P10 2015 OS) \\

\parallelenpaar{44}{}{}{}{}
α = \leerfeld °
\end{teile}
\end{aufgabe}

%% TODO Titel: Ich-kann-Satz fehlt in der Bank (Platzhalter: Kettenname) – Nr. 14
%% TODO Anweisung: ein Satz über der ersten Teilaufgabe fehlt in der Bank – Nr. 14
\begin{aufgabe}{Parallelität aus gleichen Stufenwinkeln prüfen}
\begin{teile}
\teil Die Gerade k schneidet die Geraden g (oben) und h (unten). An g liegt oberhalb von g rechts von k ein Winkel von $73^\circ$, an h oberhalb von h rechts von k ebenfalls $73^\circ$. Sind g und h parallel? \\ \kreuz{ja} \\ \kreuz{nein}
\end{teile}
\end{aufgabe}

%% TODO Titel: Ich-kann-Satz fehlt in der Bank (Platzhalter: Kettenname) – Nr. 15
%% TODO Anweisung: ein Satz über der ersten Teilaufgabe fehlt in der Bank – Nr. 15
\begin{aufgabe}{Winkel an Geraden – Fehler finden, Begründen, Anwendung}
\begin{teile}
\teil An einer Geradenkreuzung ist $\alpha = 66^\circ$. Jonas berechnet den Scheitelwinkel γ: \rechnung{\gamma &= 180^\circ - 66^\circ = 114^\circ} Finde den Fehler und gib γ richtig an.
\teil An einer Geradenkreuzung liegen α und γ einander gegenüber, β liegt dazwischen. Warum sind α und γ immer gleich groß? Begründe, ohne zu messen.
\teil Auf einem Parkplatz sind die Parkstreifen parallel. Der erste Streifen trifft den geraden Bordstein unter $62^\circ$. Welche beiden Winkel bildet der fünfte Streifen mit dem Bordstein? \leerfeld ° und \leerfeld °
\end{teile}
\end{aufgabe}

\clearpage
% Einheit 3 von 5 – Winkelsummen, Dreiecke und Vierecke (bank/winkel-dreiecke/e3.jsonl, zusammenbau v0.1)
%% TODO Zweigzeile Teil 1: was hier gelernt wird (Satz in Schülersprache aus dem Einheitstitel) – Einheit 3
\einheitenkopf[e3]{Einheit 3 von 5 · Winkelsummen, Dreiecke und Vierecke}
\zweigzeile{ab Kl. 5 · P10 oft}
\setcounter{aufgabe}{15}

%% TODO Titel: Ich-kann-Satz fehlt in der Bank (Platzhalter: Kettenname) – Nr. 16
%% TODO Anweisung: ein Satz über der ersten Teilaufgabe fehlt in der Bank – Nr. 16
\begin{aufgabe}{Welches Teildreieck?}
\begin{teile}
\teil Im Dreieck ABC ist CD die Höhe auf AB. Gesucht ist der Winkel bei C zwischen AC und CD. Fahre in einer Skizze das Dreieck nach, in dem er liegt, und kreuze an: \\ \kreuz{Dreieck ADC} \\ \kreuz{Dreieck DBC} \\ \kreuz{Dreieck ABC}
\end{teile}
\end{aufgabe}

%% TODO Titel: Ich-kann-Satz fehlt in der Bank (Platzhalter: Kettenname) – Nr. 17
%% TODO Anweisung: ein Satz über der ersten Teilaufgabe fehlt in der Bank – Nr. 17
\begin{aufgabe}{Winkelsumme}
%% TODO \rechenplatz{n}: Schrittzahl des Grundfalls fehlt in der Bank (nur bei fester Schreibform, 2.3 b)
\begin{teile}
\teil Markiere gleich lange Seiten oder den rechten Winkel und kreuze die Dreiecksart an. \\ \kreuz{spitzwinklig} \\ \kreuz{rechtwinklig} \\ \kreuz{stumpfwinklig}

\dreieckrw{4}{3}{a}{b}{c}
\teil Im Dreieck ABC ist $\alpha = 52^\circ$ und $\beta = 71^\circ$. Wie groß ist γ? γ = \leerfeld °
\teil Ein dreieckiges Grundstück hat an zwei Ecken die Winkel $51{,}4^\circ$ und $76{,}9^\circ$. Wie groß ist der Winkel an der dritten Ecke? \hfill (P10 2017 OS) \\ \leerfeld °
\teil Ein gleichschenkliges Dreieck hat zwei Basiswinkel von je $67^\circ$. Wie groß ist der Winkel an der Spitze? \leerfeld °
\teil Ein Dreieck hat drei gleich lange Seiten. Wie groß ist jeder seiner Winkel? \leerfeld °
\teil Ein Dreieck hat einen rechten Winkel und einen Winkel von $38^\circ$. Wie groß ist der dritte Winkel? \leerfeld °
\teil Ein Viereck hat die Winkel $84^\circ$, $97^\circ$ und $103^\circ$. Wie groß ist der vierte Winkel? \leerfeld °
\teil Im Giebeldreieck ABC eines Hauses steht die Stütze CD senkrecht auf dem Balken AB. An der Ecke A ist $\alpha = 27{,}6^\circ$. Wie groß ist der Winkel $\gamma_1$ zwischen AC und der Stütze CD? \hfill (P10 2022 OS) \\ γ₁ = \leerfeld °
\teil Im Dreieck ABC ist $\alpha = \beta = 64^\circ$ und AC = 5 cm. Welche Seite ist genauso lang wie AC, und wie lang ist sie?
\teil Im Dreieck ABC sind die Winkel bei A und bei B je halb so groß wie ein rechter Winkel. Begründe, dass der Winkel bei C ein rechter ist.
\teil Ergänze richtig: „In jedem Drachenviereck …“ \hfill (P10 2026 FOR) \\ \kreuz{sind alle Winkel gleich groß} \\ \kreuz{gibt es zwei Paare gleich langer Nachbarseiten} \\ \kreuz{sind alle Seiten gleich lang}
\teil Im Dreieck ABC ist $\alpha = \beta = 66^\circ$ und AC = 7 cm. Gib die Länge von BC und die Größe von γ an. \hfill (P10 2023 OS) \\ BC = \leerfeld[cm] , γ = \leerfeld °
\end{teile}
\end{aufgabe}

%% TODO Titel: Ich-kann-Satz fehlt in der Bank (Platzhalter: Kettenname) – Nr. 18
%% TODO Anweisung: ein Satz über der ersten Teilaufgabe fehlt in der Bank – Nr. 18
\begin{aufgabe}{Innenwinkelsumme von Fünf- und Sechseck über Dreiecke}
\begin{teile}
\teil Zerlege ein Fünfeck durch Diagonalen von einer Ecke aus in Dreiecke. Wie groß ist die Summe seiner Innenwinkel? \leerfeld °
\end{teile}
\end{aufgabe}

%% TODO Titel: Ich-kann-Satz fehlt in der Bank (Platzhalter: Kettenname) – Nr. 19
%% TODO Anweisung: ein Satz über der ersten Teilaufgabe fehlt in der Bank – Nr. 19
\begin{aufgabe}{Winkelsumme – Fehler finden, Begründen, Anwendung}
\begin{teile}
\teil Ein Viereck hat die Winkel $52^\circ$, $71^\circ$ und $93^\circ$. Ben rechnet den vierten Winkel: \rechnung{52^\circ + 71^\circ + 93^\circ &= 216^\circ \\ 216^\circ - 180^\circ &= 36^\circ} Finde den Fehler und rechne richtig.
\teil Zeichne durch C die Parallele zu AB. Begründe mit Wechselwinkeln, dass α, β und γ zusammen einen gestreckten Winkel ergeben.

\dreieck{(0,0)}{(5,0)}{(1.8,3)}{a}{b}{c}{\alpha}{\beta}{\gamma}
\teil Ein Satteldach ist im Querschnitt ein gleichschenkliges Dreieck; beide Dachflächen sind um $38^\circ$ gegen die Waagerechte geneigt. Wie groß ist der Winkel am First? \leerfeld °
\end{teile}
\end{aufgabe}

\clearpage
% Einheit 4 von 5 – Dreiecke konstruieren (bank/winkel-dreiecke/e4.jsonl, zusammenbau v0.1)
%% TODO Zweigzeile Teil 1: was hier gelernt wird (Satz in Schülersprache aus dem Einheitstitel) – Einheit 4
\einheitenkopf[e4]{Einheit 4 von 5 · Dreiecke konstruieren}
\zweigzeile{ab Kl. 6, je nach Buch bis Kl. 7 (am Gymnasium ab Kl. 5) · P10}
\setcounter{aufgabe}{19}

%% TODO Titel: Ich-kann-Satz fehlt in der Bank (Platzhalter: Kettenname) – Nr. 20
%% TODO Anweisung: ein Satz über der ersten Teilaufgabe fehlt in der Bank – Nr. 20
\begin{aufgabe}{Konstruieren}
%% TODO \rechenplatz{n}: Schrittzahl des Grundfalls fehlt in der Bank (nur bei fester Schreibform, 2.3 b)
\begin{teile}
\teil Gegeben sind a = 6 cm, b = 4 cm, $\gamma = 52^\circ$. Markiere die Stücke in einer Planfigur. Welcher Kongruenzsatz passt? \\ \kreuz{SSS} \\ \kreuz{SWS} \\ \kreuz{WSW} \\ \kreuz{SsW}
\teil Konstruiere das Dreieck ABC mit a = 6 cm, b = 5 cm und c = 7 cm. Beginne mit der Seite c auf dem Strahl.

\winkelstrahl[9][A]
\teil Konstruiere das Dreieck ABC mit b = 5 cm, c = 6 cm und $\alpha = 50^\circ$. Beginne mit der Seite c auf dem Strahl.

\winkelstrahl[9][A]
\teil Konstruiere das Dreieck ABC mit c = 7 cm, $\alpha = 38^\circ$ und $\beta = 71^\circ$. Beginne mit der Seite c auf dem Strahl.

\winkelstrahl[9][A]
\teil Konstruiere das Dreieck ABC mit b = 6 cm, c = 5 cm und $\beta = 64^\circ$. Beginne mit der Seite c auf dem Strahl.

\winkelstrahl[9][A]
\teil Konstruiere ein gleichschenkliges Dreieck ABC mit der Basis c = 6 cm und den Schenkeln a = b = 5 cm.

\winkelstrahl[9][A]
\teil Konstruiere ein rechtwinkliges Dreieck ABC mit dem rechten Winkel bei C und den Katheten a = 4 cm und b = 6 cm.

\winkelstrahl[9][C]
\teil Schreibe eine Konstruktionsbeschreibung in drei bis vier Schritten für das Dreieck ABC mit c = 6 cm, $\alpha = 55^\circ$ und b = 4 cm.
\teil Lässt sich aus den Seiten 4 cm, 5 cm und 10 cm ein Dreieck konstruieren? \\ \kreuz{ja} \\ \kreuz{nein}
\teil Dreieck ABC: AB = 5 cm, BC = 7 cm, CA = 6 cm. Dreieck PQR: PQ = 7 cm, QR = 6 cm, RP = 5 cm. Die Dreiecke sind kongruent. Welche Ecke von PQR gehört zu A, zu B und zu C?
\teil Ein dreieckiges Beet ABC hat die Seite a = 12 m. Was gilt für die Summe der beiden anderen Seiten b und c? \hfill (P10 2020 OS) \\ \kreuz{$b + c > 12$ m} \\ \kreuz{$b + c = 12$ m} \\ \kreuz{$b + c < 12$ m}
\end{teile}
\end{aufgabe}

%% TODO Titel: Ich-kann-Satz fehlt in der Bank (Platzhalter: Kettenname) – Nr. 21
%% TODO Anweisung: ein Satz über der ersten Teilaufgabe fehlt in der Bank – Nr. 21
\begin{aufgabe}{Dreieck skizzieren und beschriften}
\begin{teile}
\teil Beschrifte die Planfigur: Ecken A, B, C gegen den Uhrzeigersinn, dazu die Seiten a, b, c und die Winkel α, β, γ. Markiere farbig die gegebenen Stücke b = 5 cm, c = 6 cm und $\alpha = 38^\circ$.

\dreieck{(0,0)}{(5,0)}{(1.5,3)}{}{}{}{}{}{}
\end{teile}
\end{aufgabe}

%% TODO Titel: Ich-kann-Satz fehlt in der Bank (Platzhalter: Kettenname) – Nr. 22
%% TODO Anweisung: ein Satz über der ersten Teilaufgabe fehlt in der Bank – Nr. 22
\begin{aufgabe}{Konstruieren – Fehler finden, Begründen, Anwendung}
\begin{teile}
\teil Sina sagt: „Mit $\alpha = 52^\circ$, $\beta = 64^\circ$ und $\gamma = 64^\circ$ ist das Dreieck eindeutig festgelegt – das ist der Kongruenzsatz WWW.“ Finde den Fehler.
\teil Warum reicht es nicht, von einem Dreieck nur die drei Winkel zu kennen? Begründe.
\teil Ein dreieckiges Beet soll mit Brettern von 3,5 m, 2,0 m und 1,2 m Länge eingefasst werden. Geht das? Begründe.
\end{teile}
\end{aufgabe}

\clearpage
% Einheit 5 von 5 – Besondere Linien im Dreieck und Satz des Thales (bank/winkel-dreiecke/e5.jsonl, zusammenbau v0.1)
%% TODO Zweigzeile Teil 1: was hier gelernt wird (Satz in Schülersprache aus dem Einheitstitel) – Einheit 5
\einheitenkopf[e5]{Einheit 5 von 5 · Besondere Linien im Dreieck und Satz des Thales}
\zweigzeile{ab Kl. 6, je nach Buch bis Kl. 7 · P10}
\setcounter{aufgabe}{22}

%% TODO Titel: Ich-kann-Satz fehlt in der Bank (Platzhalter: Kettenname) – Nr. 23
%% TODO Anweisung: ein Satz über der ersten Teilaufgabe fehlt in der Bank – Nr. 23
\begin{aufgabe}{Besondere Linien}
%% TODO \rechenplatz{n}: Schrittzahl des Grundfalls fehlt in der Bank (nur bei fester Schreibform, 2.3 b)
\begin{teile}
\teil Eine Linie geht durch die Ecke C und steht senkrecht auf der Seite AB. Welche Linie im Dreieck ABC ist das? \\ \kreuz{Mittelsenkrechte} \\ \kreuz{Winkelhalbierende} \\ \kreuz{Höhe} \\ \kreuz{Seitenhalbierende}
\teil Verbinde A und B und konstruiere mit dem Zirkel die Mittelsenkrechte der Strecke AB. In welchem Punkt schneidet sie AB?

\begin{ksys}[xmin=0,xmax=10,ymin=0,ymax=6] \punkt{2}{2}{A} \punkt{8}{4}{B} \end{ksys}
M( \leerfeld | \leerfeld )
\teil Zeichne das spitzwinklige Dreieck ABC. Konstruiere zwei Mittelsenkrechten und den Umkreis. Gib den Mittelpunkt U an.

\begin{ksys}[xmin=0,xmax=10,ymin=0,ymax=8] \punkt{1}{1}{A} \punkt{7}{1}{B} \punkt{3}{5}{C} \end{ksys}
U( \leerfeld | \leerfeld )
\teil Konstruiere mit dem Zirkel die Winkelhalbierende von α. Wie groß ist jeder der beiden Teilwinkel?

\winkel[5]{64}{\alpha}
\leerfeld °
\teil Zeichne das Dreieck ABC. Konstruiere zwei Winkelhalbierende und den Inkreis. Gib den Mittelpunkt I an.

\begin{ksys}[xmin=0,xmax=10,ymin=0,ymax=10] \punkt{1}{1}{A} \punkt{7}{1}{B} \punkt{1}{9}{C} \end{ksys}
I( \leerfeld | \leerfeld )
\teil Zeichne das Dreieck ABC und mit dem Geodreieck die Höhe von C auf die Seite AB, wenn nötig auf ihre Verlängerung. Gib den Fußpunkt D an.

\begin{ksys}[xmin=0,xmax=10,ymin=0,ymax=8] \punkt{1}{1}{A} \punkt{7}{1}{B} \punkt{4}{5}{C} \end{ksys}
D( \leerfeld | \leerfeld )
\teil Zeichne das Dreieck ABC und seine drei Seitenhalbierenden. Gib den Schwerpunkt S an.

\begin{ksys}[xmin=0,xmax=10,ymin=0,ymax=8] \punkt{1}{1}{A} \punkt{7}{1}{B} \punkt{4}{7}{C} \end{ksys}
S( \leerfeld | \leerfeld )
\teil Zeichne eine Strecke AB von 8 cm Länge und darüber den Thaleskreis. Wähle einen Punkt C darauf und zeichne das Dreieck ABC. Welchen Radius hat der Halbkreis, und wo liegt der rechte Winkel?

\winkelstrahl[10][A]
r = \leerfeld[cm]
\teil Die Punkte A und B liegen sich auf einem Kreis mit dem Mittelpunkt M gegenüber, die Strecke AB geht durch M. C ist ein weiterer Punkt auf dem Kreis. Begründe, dass das Dreieck ABC bei C einen rechten Winkel hat.
\teil Im Viereck ABCD schneiden sich die Diagonalen AC und BD im Punkt M. Es gilt MA = MC = MD = 4 cm. Begründe, dass der Winkel des Vierecks bei D ein rechter ist.
\end{teile}
\end{aufgabe}

%% TODO Titel: Ich-kann-Satz fehlt in der Bank (Platzhalter: Kettenname) – Nr. 24
%% TODO Anweisung: ein Satz über der ersten Teilaufgabe fehlt in der Bank – Nr. 24
\begin{aufgabe}{Besondere Linien – Fehler finden, Begründen, Anwendung}
\begin{teile}
\teil Lukas soll im Dreieck ABC die Höhe von C auf AB zeichnen. Er verbindet C mit der Mitte von AB. Finde den Fehler.
\teil Warum ist der Schnittpunkt der Mittelsenkrechten eines Dreiecks ABC von allen drei Ecken gleich weit entfernt? Begründe.
\teil Drei Dörfer A, B und C wollen einen gemeinsamen Funkmast, der von allen drei Dörfern gleich weit entfernt ist (1 Kästchen = 1 km). Konstruiere den Standort. Gib seine Koordinaten und die Entfernung zu den Dörfern an.

\begin{ksys}[xmin=0,xmax=10,ymin=0,ymax=8] \punkt{1}{1}{A} \punkt{9}{1}{B} \punkt{1}{7}{C} \end{ksys}
( \leerfeld | \leerfeld )
\end{teile}
\end{aufgabe}

% Abhakseite „Das kann ich“ – Zone nur im Gesamt (\mitzone)
%% TODO Ich-kann-Sätze: Platzhalter wie in den Aufgabendateien
\begin{abhakseite}
\ifdefined\mitzone
\abhakgruppe{Kennst du schon}
\abhak{1}{Mit Lineal und Geodreieck zeichnen}
\abhak{2}{Addieren und Subtrahieren im Bereich bis zum Vollwinkel, auch mit Dezimalzahlen (Vielfaches minus zwei Dezimalgrade)}
\abhak{3}{Längeneinheiten m und km umrechnen (1,5 km = 1500 m)}
\abhak{4}{Vierecksarten kennen}
\abhak{5}{Kreis mit dem Zirkel zeichnen, Radius und Durchmesser unterscheiden}
\abhak{6}{Addieren und Subtrahieren im Bereich bis zum Vollwinkel, auch mit Dezimalzahlen (Vielfaches minus zwei Dezimalgrade) – Fehler finden}
\abhak{7}{Addieren und Subtrahieren im Bereich bis zum Vollwinkel, auch mit Dezimalzahlen (Vielfaches minus zwei Dezimalgrade) – selbst rechnen}
\fi
\abhakgruppe{Einheit 1 · Winkel messen und zeichnen}
\abhak{8}{Welche Skala?}
\abhak{9}{Winkel messen}
\abhak{10}{Winkel schätzen}
\abhak{11}{Winkel messen – Fehler finden, Begründen, Anwendung}
\abhakgruppe{Einheit 2 · Winkel an Geradenkreuzungen und Parallelen}
\abhak{12}{Sind die Geraden parallel?}
\abhak{13}{Winkel an Geraden}
\abhak{14}{Parallelität aus gleichen Stufenwinkeln prüfen}
\abhak{15}{Winkel an Geraden – Fehler finden, Begründen, Anwendung}
\abhakgruppe{Einheit 3 · Winkelsummen, Dreiecke und Vierecke}
\abhak{16}{Welches Teildreieck?}
\abhak{17}{Winkelsumme}
\abhak{18}{Innenwinkelsumme von Fünf- und Sechseck über Dreiecke}
\abhak{19}{Winkelsumme – Fehler finden, Begründen, Anwendung}
\abhakgruppe{Einheit 4 · Dreiecke konstruieren}
\abhak{20}{Konstruieren}
\abhak{21}{Dreieck skizzieren und beschriften}
\abhak{22}{Konstruieren – Fehler finden, Begründen, Anwendung}
\abhakgruppe{Einheit 5 · Besondere Linien im Dreieck und Satz des Thales}
\abhak{23}{Besondere Linien}
\abhak{24}{Besondere Linien – Fehler finden, Begründen, Anwendung}
\end{abhakseite}

\end{document}
